\chapter{第二十四章：变革联盟与软着陆路径}

# 第二十四章：变革联盟与软着陆路径

\textit{"改革如果没有同盟军，就只能孤军奋战。"}
——俾斯麦

\hline\newpage
---

老王我当年在工厂推行过一个改革。不是什么大事，就是把车间里的物料堆放位置重新调整了一下，按理说应该提高效率。

结果呢？物料部门说我们物料多不方便挪；生产部门说换位置影响他们干活；后勤部门说调整得给钱；财务部门说这算不算资本性支出要不要审批。

我说："就挪几个架子，不用这么多程序吧？"

他们说："你不懂，这牵一发动全身。"

后来我才明白：任何变革，表面上是改事，实际上是改人。任何改革，表面上是业务流程，实际上是利益格局。

光靠你一个人喊"我们要改革"，没人响应，最后只能你自己累个半死还落一身埋怨。所以变革这事，得有同盟军。

\section{一、变革联盟的必要性}

## 一、变革联盟的必要性

任何成功的组织变革，都需要一个变革联盟——一群有足够影响力的人，愿意并且能够推动变革向前进。

这个联盟的作用是提供变革的政治支持。当变革遇到阻力时，变革联盟可以帮助克服。当变革遇到资源竞争时，变革联盟可以帮助争取。

在AI变革中，变革联盟尤为重要，因为AI变革几乎必然触动既得利益者。没有足够强大的变革联盟，AI变革很容易在既得利益者的阻力下停滞。

但建立一个有效的变革联盟，比听起来要困难得多。

就像你组织一场饭局，你得请那些"能说话算话的人"，还得请那些"能带动气氛的人"，还得请那些"能买单的人"。请错了人，饭局就黄了。

\section{二、如何建立变革联盟}

## 二、如何建立变革联盟

\textbf{第一步：识别关键位置的人}

变革联盟的核心，应该包括那些在组织中占据关键位置的人——不只是正式的管理者，也包括那些对其他人有影响力的非正式领导者。

关键位置包括：那些与多个部门有联系的人（他们知道很多信息），那些在员工中有高度信任的人（他们的话有分量），那些掌握关键资源的人（他们可以提供资源）。

\textbf{第二步：找到共同愿景}

变革联盟的成员可能有不同的动机和利益。但他们必须有一个共同的愿景——他们为什么聚在一起做这件事。

这个共同愿景，是维持变革联盟凝聚力的粘合剂。

就像你打游戏组队，你得先确定你们的目标是什么——是推塔还是刷野还是拿龙。目标不一致，队伍就散了。

\textbf{第三步：让联盟成员有"拥有感"}

如果变革联盟的成员只是被"要求"支持变革，而不是真正"想要"支持变革，那这个联盟是脆弱的。

让联盟成员有拥有感，意味着让他们参与变革的设计和决策，而不只是执行命令。

\section{三、软着陆路径的设计原则}

## 三、软着陆路径的设计原则

当变革联盟建立之后，下一个问题是如何设计变革的路径，让它能够"软着陆"——让阻力最小化，让支持最大化。

\textbf{原则一：小步快跑，而非大步推进}

大的、激进的变革，往往会遭遇强大的阻力。相比之下，小步快跑的渐进式变革，更容易被接受。

AI变革的具体实践可以是：先在局部范围进行试点，证明效果后再逐步扩大，而不是一开始就在全公司推行。

\textbf{原则二：先易后难，而非一开始就碰硬骨头}

在变革的早期，建立"速赢"很重要——那些能够快速实现、容易看到效果的变革。

这些"速赢"可以增强组织对变革的信心，为面对更大的挑战积累势能。

\textbf{原则三：持续沟通，而非一次性宣告}

变革不是一次性的宣告，而是一个持续的过程。

持续沟通意味着：定期更新变革进展，及时回应员工的担忧和疑问，根据反馈调整变革方案。

\textbf{原则四：培养变革大使，而非依赖顶层推动}

如果变革只能依赖高层的推动，那一旦高层注意力转移，变革就会停滞。

培养变革大使——那些在各个层级、各个部门积极推动变革的人——可以让变革更加持续。

\section{四、如何处理阻力}

## 四、如何处理阻力

即使有了变革联盟和软着陆路径，变革过程中仍然会遇到阻力。关键是学会处理阻力，而不是被阻力击败。

\textbf{方法一：区分理性阻力和情感阻力}

阻力有时来自理性分析——变革确实存在问题，需要被解决。但阻力有时来自情感——恐惧、不确定性、被威胁感。

处理理性阻力需要数据和论证。处理情感阻力需要同理心、倾听和重新连接。

很多时候，阻力的表层是理性分析，深层是情感恐惧。只处理理性层面而不触及情感层面，阻力会持续存在。

\textbf{方法二：将阻力者转化为支持者}

那些最强烈反对变革的人，往往也是最有影响力的人。如果能够把他们转化为支持者，他们会从阻力和推动。

转化的方法通常不是反驳他们的论点，而是理解他们阻力的真正原因，针对那个原因重新设计变革方案或者提供额外的支持。

\textbf{方法三：接受"不是所有人都会支持"}

在任何变革中，总有一些人不会支持。这个事实需要被接受。

试图说服所有人支持变革，往往是浪费资源。更有效的方法是：识别那些可能支持的人，争取他们的支持；对于坚定反对者，确保他们不会破坏变革；对于中间者，持续地沟通和影响。

\section{五、衡量变革进展}

## 五、衡量变革进展

没有衡量的变革，容易失去方向和焦点。

在AI变革中，建立清晰的衡量指标，可以帮助你追踪进展、识别问题、调整方向。

衡量AI变革进展的指标可以包括：

\textbf{采用率指标：} AI工具的实际使用率、使用频率、使用深度。

\textbf{能力建设指标：} 完成AI培训的员工比例、员工AI能力评估分数、AI技能认证数量。

\textbf{业务影响指标：} AI带来的效率提升、错误率降低、新产品或新服务数量。

\textbf{文化变革指标：} 员工对AI变革的态度调查、管理者对AI的支持程度、跨部门协作频率。

\section{六、结语：变革是旅程，不是项目}

## 六、结语：变革是旅程，不是项目

变革管理中有一个常见的错误：把变革当作一个项目来做——有明确的项目启动日期，有明确的结束目标，有明确的项目团队。

但真正的变革，不是项目，而是旅程。

旅程的特点是：没有固定的终点，需要持续的适应，没有可以放下的责任。

AI变革同样如此。它不是一个可以"完成"然后"继续前进"的项目，而是一个需要组织持续投入、持续学习、持续适应的过程。

那些把变革当作旅程来管理的组织，比那些把变革当作项目来管理的组织，更有可能实现真正的转型。

俾斯麦说"改革如果没有同盟军，就只能孤军奋战"——这话是真的。但光有同盟军不够，你还得有正确的行军路线，不然同盟军越多，混乱越多。

\hline\newpage
---

\textbf{本章小结：}

1. 变革联盟是变革成功的必要条件——提供政治支持，克服阻力，争取资源
2. 建立变革联盟的步骤：识别关键位置的人、找到共同愿景、让成员有拥有感
3. 软着陆路径的设计原则：小步快跑，先易后难、持续沟通，培养变革大使
4. 处理阻力的方法：区分理性阻力和情感阻力、将阻力者转化为支持者、接受不是所有人都会支持
5. 衡量变革进展：采用率、能力建设、业务影响，文化变革四类指标
6. 变革是旅程不是项目——需要持续投入、持续学习、持续适应
